% Source: https://artofproblemsolving.com/community/c3h2256174_a_pretty_cool_formula_i_discovered
% Stable identity: aops-topic-2256174-q01
\textbf{(a)} Let $S=1+2+\cdots+n$. Writing the sum forwards and backwards and adding term by term gives $n$ pairs, each with sum $n+1$, so $2S=n(n+1)$ and $S=\dfrac{n(n+1)}{2}$. Factoring out $d$,
\[d+2d+\cdots+nd=d(1+2+\cdots+n)=\dfrac{dn(n+1)}{2}.\]

\textbf{(b)} Here $d=7$ and $nd=700$, so $n=100$:
\[7+14+\cdots+700=\dfrac{7\cdot100\cdot101}{2}=\boxed{35350}.\]

\textbf{(c)} The sum consists of $n$ copies of $a$ plus $d\bigl(0+1+\cdots+(n-1)\bigr)$, so
\[a+(a+d)+\cdots+\bigl(a+(n-1)d\bigr)=na+\dfrac{n(n-1)}{2}\,d=\boxed{\dfrac{n\bigl(2a+(n-1)d\bigr)}{2}},\]
that is, $n$ times the average of the first and last terms. Part (a) is the case $a=d$.
